\clearpage
\VerdierPhysicalPageAnchor{0122}%
\begingroup
\small
\setlength{\emergencystretch}{1.5em}%
\setlength{\abovedisplayskip}{4pt plus 2pt minus 2pt}%
\setlength{\belowdisplayskip}{4pt plus 2pt minus 2pt}%

\VerdierSemanticAnchor{ch2-1-2-13-countability-hypothesis}
\textit{Si les produits dénombrables ou les sommes directes dénombrables
sont représentables dans $\mathcal D$, la catégorie $\mathcal D$ est
décomposable.}

\VerdierSemanticAnchor{ch2-1-2-13-proof-morphism}
En effet, soit $v\colon W\longrightarrow X$ un morphisme tel que $fv=0$.
On a alors un morphisme de morphismes :

\VerdierSemanticAnchor{ch2-1-2-13-morphism-square}
\[
  \begin{tikzcd}[column sep=4em,row sep=2.5em]
    X \arrow[r,"f"] & Y \\
    W \arrow[u,"v"] \arrow[r] & 0 \arrow[u]
  \end{tikzcd}
  \,.
\]

\VerdierSemanticAnchor{ch2-1-2-13-1}
D'où, en appliquant le foncteur $\operatorname{Tr}$, un diagramme
commutatif :
\begin{equation}
  \resizebox{0.82\linewidth}{!}{%
    $\begin{tikzcd}[ampersand replacement=\&,column sep=2.05em,row sep=3em]
      X \arrow[r,"f"]
        \& Y \arrow[r,"{a(f)}"]
        \& \tau(f) \arrow[r,"{b(f)}"]
        \& X[1] \\
      W \arrow[u,"v"] \arrow[r]
        \& 0 \arrow[u] \arrow[r]
        \& \tau(W\longrightarrow 0) \arrow[u,"{\operatorname{Tr}(v,0)}"]
          \arrow[r,"{b(W\longrightarrow 0)}"']
        \& W[1] \arrow[u,"{v[1]}"']
    \end{tikzcd}$%
  }\,.
  \tag{1.2.13.1}
\end{equation}

\VerdierSemanticAnchor{ch2-1-2-13-bottom-triangle}
Remarquons tout d'abord que la ligne du bas du diagramme (1.2.13.1) est
un triangle distingué. Par suite, le morphisme
$b(W\longrightarrow 0)$ est un isomorphisme (1.2.1). Posons alors :

\VerdierSemanticAnchor{ch2-1-2-13-alpha-definition}
\[
  \alpha(v)=\operatorname{Tr}(v,0)[-1]
  b(W\longrightarrow 0)^{-1}[-1]\,,
\]
de sorte que le diagramme :

\VerdierSemanticAnchor{ch2-1-2-13-alpha-diagram}
\[
  \begin{tikzcd}[column sep=3.5em,row sep=2.4em]
    \tau(f)[-1] \arrow[r,"{b(f)[-1]}"]
      & X \arrow[r,"f"]
      & Y \\
      & W \arrow[ul,"{\alpha(v)}"'] \arrow[u,"v"] \arrow[ur,"0"']
  \end{tikzcd}
\]
soit commutatif. On vérifie aussitôt qu'on a ainsi défini une section :
\par

\VerdierSemanticAnchor{ch2-1-2-13-alpha-section}
\[
  \alpha\colon \operatorname{Ker}(f)\longrightarrow \tau(f)[-1]
\]
\noindent du morphisme canonique
$\tau(f)[-1]\longrightarrow\operatorname{Ker}(f)$ du foncteur représenté
par $\tau(f)[-1]$ dans le foncteur noyau du morphisme $f$.

\VerdierSemanticAnchor{ch2-1-2-13-projector-kernel}
Le foncteur noyau de $f$ est donc isomorphe au foncteur noyau d'un
projecteur de $\tau(f)[-1]$. Il est donc représentable (1.2.9).

\VerdierSemanticAnchor{ch2-1-2-13-proof-complete}
On en déduit par (1.2.9) que le morphisme $f$ est décomposable, ce qui
achève la démonstration.
\endgroup
